% Figaro Paper K: A Permissionless, Ownerless Container-Shipping Assembly
% Industrial / Systems Engineering register; the TradeLens replacement
% worked example.
% Audience: industrial-engineering and systems-engineering scholars,
% supply-chain-coordination researchers, maritime-shipping practitioners,
% trade-finance professionals.
% Companion to Paper J (scheduled-service / Spirit Air worked example).
% Preprint, May 2026

\documentclass[11pt,a4paper]{article}

\usepackage[utf8]{inputenc}
\usepackage[T1]{fontenc}
\usepackage{amsmath,amssymb,amsthm}
\usepackage{mathtools}
\usepackage{booktabs}
\usepackage{graphicx}
\usepackage{hyperref}
\usepackage{cleveref}
\usepackage{enumitem}
\usepackage{xcolor}
\usepackage[margin=1in]{geometry}
\usepackage{natbib}
\bibliographystyle{plainnat}

\newtheorem{theorem}{Theorem}[section]
\newtheorem{proposition}[theorem]{Proposition}
\newtheorem{corollary}[theorem]{Corollary}
\theoremstyle{definition}
\newtheorem{definition}[theorem]{Definition}
\newtheorem{example}[theorem]{Example}
\theoremstyle{remark}
\newtheorem{remark}[theorem]{Remark}

\newcommand{\figaro}{\textsc{Figaro}}
\newcommand{\pay}{P}
\newcommand{\cumval}{G}

\title{%
  A Permissionless, Ownerless Supply-Chain Assembly:\\
  Why TradeLens Failed and the Consumer-Rooted Architecture That Succeeds It%
}

\author{
  Alessandro Daliana\thanks{Corresponding author. \texttt{adaliana@figaro.org}}
}

\date{May 2026}

\begin{document}

\maketitle

% ════════════════════════════════════════════════════════════════════
\begin{abstract}
TradeLens (2018--2023), the IBM--Maersk container-shipping consortium
that aimed to produce a tamper-proof shared record of every container's
journey, shut down with only one major non-Maersk carrier on board. The
shutdown was not a software failure. It was a governance failure
visible from the consortium's structural shape: a competitor-controlled
platform, requiring participation from carriers who would not ratify a
competitor's gatekeeping over an industry resource, with high
integration cost and no settlement guarantee binding parties to the
arrangement. We argue the failure was structural, predictable from the
shape of the architecture, and not specific to TradeLens.

TradeLens addressed the maritime portion of a longer chain. The
architectural alternative we develop reaches further: from factory
to consumer. We present a permissionless, ownerless supply-chain
assembly built on the bonded settlement primitive of the companion
mechanism paper. The assembly has no consortium, no foundation, no
admin function, no upgrade path, and no central party who can be
ratified or refused. The \emph{final consumer} is the
\texttt{rootBuyer} of the supply-chain process, committing directly
to every party that contributes to the goods arriving at the
consumer's door --- factory, packaging, origin logistics, port
authorities at origin and destination, ocean carrier, customs
broker, customs authority (for import duties), domestic logistics,
state government (for sales tax), and any other resource provider
the goods draw on. Carriers who want to participate compose;
carriers who do not do not, with no consequence for either decision.
Bonded commitments at each commit lock skin in the game proportional
to the cumulative process value; atomic resolution by the consumer
settles the whole process on receipt of the goods. The firm-shape
intermediaries the conventional supply chain inserts at every tier
(shipper-of-record, distributor, wholesaler, retailer) become
structurally unnecessary because the consumer's bonded commits reach
the operational sellers directly.
The schemas the assembly draws on are largely existing infrastructure
(commerce, geo, handoff, jurisdiction, fulfilment, GHG family,
proximity); two new schemas (container-seal, Incoterms-2020-as-anchor)
extend the protocol's surface in ways that respect the
extension-doctrine. We also treat the bill-of-lading transferability
question through the cancellable-seller / counter-process pattern that
expresses transferability at the protocol layer without weakening the
kernel's invariants.

The paper is in industrial-engineering register: the contribution is
the assembly DAG and its operational properties, not novel
mechanism-design results. Why TradeLens failed and what its successor
looks like architecturally are answerable as supply-chain coordination
questions.
\end{abstract}

\medskip
\noindent\textbf{Keywords:}
container shipping, supply-chain coordination, bills of lading,
Incoterms, MLETR, transferable records, process modeling

% ════════════════════════════════════════════════════════════════════
\section{Introduction}
\label{sec:intro}

Container shipping is a coordination problem of unusual scale: a
single 20{,}000-TEU vessel carries cargo from thousands of shippers,
moves through multiple ports under multiple jurisdictions, transfers
to multiple consignees, and is documented by a stack of bills of
lading, customs filings, GHG attestations, and operational handoffs
that no single party originates and no single party owns. The
information problem is real --- duplicated paper-based records,
opaque vessel-position data, contested seal-integrity claims --- and
the industry has been searching for a digital solution for two
decades.

TradeLens, the consortium platform jointly developed by IBM and
Maersk between 2018 and 2023, was the most ambitious attempt at
that solution. Its approach was a permissioned blockchain on which
participants would publish operational records that all parties to
a shipment could read. The platform shut down in early 2023 after
failing to attract more than one major non-Maersk carrier (CMA CGM)
into participation; MSC, Hapag-Lloyd, ONE, and Evergreen all
declined. The platform was technically functional; the consortium
was not.

\paragraph{Reading the failure structurally.}
The standard reading of TradeLens's failure focuses on adoption
dynamics: the network-effects argument that any shared-ledger
platform requires a critical-mass coalition before it is useful.
This is not wrong, but it is incomplete. The critical-mass argument
explains why TradeLens \emph{failed to grow}; it does not explain
why the carriers who declined to join made what was, from their
perspective, a rational decision. The carriers' decision was not
about technology adoption. It was about ratifying a competitor's
control of an industry resource, which a competitor-controlled
platform necessarily asks them to do.

The \emph{governance} structure of the platform was the binding
constraint, not the network effects. A competitor-controlled
platform requires the other competitors to accept the gatekeeping
authority of the controlling competitor; the other competitors
have correctly read this as conceding strategic position and have
declined. The network-effects collapse follows from the governance
choice; it does not cause it. A platform with neutral governance
might have faced its own difficulties, but it would not have faced
the specific difficulty TradeLens did.

\paragraph{The architectural alternative.}
We argue that the structural fix is not better governance but no
governance at the platform layer at all. A permissionless,
ownerless settlement primitive --- the bonded commitment of the
companion mechanism paper --- removes the governance question from
the architecture. There is no consortium to join; there is a
public protocol that any wallet can compose with. There is no
gatekeeping authority to refuse; carriers participate by signing
commitments and decline by not signing. The network-effects
question becomes more manageable because the platform-as-veto
problem disappears: a competitor's participation does not require
the competitor's consent to a competitor's gatekeeping.

The assembly described in this paper is what supply-chain
coordination looks like under such a substrate. TradeLens
addressed the maritime portion of the chain --- carrier-to-carrier
and carrier-to-port records of cargo movement; the bonded
architecture reaches further. It admits the \emph{full chain from
factory to consumer} under a single settlement primitive, with
the final consumer as the \texttt{rootBuyer} of the supply-chain
process and every operational party (factory, packaging, origin
logistics, port authorities, ocean carrier, customs broker,
customs authority for tariffs, domestic logistics, state
government for sales tax, and any other resource provider the
goods draw on) as a seller in that process. The firm-shape
intermediaries of the conventional chain --- shipper-of-record,
distributor, wholesaler, retailer --- become structurally
unnecessary because the consumer's bonded commits reach the
operational sellers directly. It is not a TradeLens replacement
in the sense of a drop-in alternative product; it is a
re-architecture of the supply-chain coordination problem in which
the firm-tiered chain is replaced by a consumer-rooted bonded
process.

\paragraph{Paper organization.}
\Cref{sec:tradelens} reviews the TradeLens history and the
structural reading of its failure (the maritime case).
\Cref{sec:primitive} summarizes the settlement primitive in the
register the present paper uses.
\Cref{sec:assembly} develops the consumer-rooted assembly: the
organizational DAG (factory to consumer), the schema-attested
operational handoffs along it, and the kernel commit structure
that backs it (consumer as \texttt{rootBuyer}, every commit
consumer-to-seller, all in one process).
\Cref{sec:bonds} works out the bond posture and cohort-compensation
mechanics across the resource providers, with a worked example
including realistic numbers, tariffs, and sales tax.
\Cref{sec:schemas} treats the schema requirements: most existing
infrastructure, two new schemas (container-seal-v1 and
Incoterms-2020-as-anchor), and the deliberate non-inclusion of
transferable-BoL semantics in the kernel-layer schema set.
\Cref{sec:bol} treats bill-of-lading transferability as a
protocol-layer composition pattern (CancellableSeller wrapper +
counter-process).
\Cref{sec:comparison} compares the architecture to CargoX,
TradeTrust, MLETR, and TradeLens itself.
\Cref{sec:scope} states the scope conditions and what the
architecture deliberately does not address.
\Cref{sec:conclusion} concludes.

% ════════════════════════════════════════════════════════════════════
\section{What TradeLens Was and Why It Failed}
\label{sec:tradelens}

We treat the failure structurally rather than narratively.

\paragraph{What TradeLens was.}
TradeLens was a permissioned blockchain platform for global
container shipping, jointly developed by IBM and Maersk and
operational from 2018 to 2023. The platform's stated goal was a
shared, tamper-proof record of every container's journey ---
bills of lading, customs filings, port handoffs, vessel positions
--- accessible to all parties in a shipment.

The technical architecture was Hyperledger-Fabric-based: a
permissioned ledger where participating carriers, ports, and
customs authorities could publish operational records. The
platform's value proposition rested on the assumption that all
major ocean carriers would join, producing a single ledger that
shippers could rely on as comprehensive cross-carrier reference.
\citet{jensen_etal_2019} and \citet{sarker_etal_2021} document
the platform's value proposition and use-case design as it stood
mid-deployment; both papers predate the shutdown and frame the
platform around adoption potential rather than around the
governance refusal we read as the binding constraint below.
The governance-was-binding reading is original to this paper;
Jensen and Sarker supply the value-proposition baseline against
which the reading is offered.

\paragraph{The adoption shape.}
By 2022, only CMA CGM (a Maersk competitor) had joined as a major
non-Maersk carrier. MSC (the largest carrier by capacity),
Hapag-Lloyd, Ocean Network Express (ONE), and Evergreen had all
declined to participate. In November 2022, IBM and Maersk
announced the platform's shutdown; production ended in Q1 2023
\citep{lloyd_loadstar_2022}.

\paragraph{Why the carriers declined.}
The standard adoption-dynamics reading explains the shutdown
through a familiar mechanism: shippers want a single ledger
covering 90\%+ of their cargo volume; without MSC and Hapag-Lloyd
the data was incomplete; without complete data, shippers wouldn't
pay for the platform; without revenue, carriers had no
incremental reason to join; the equilibrium was reflexive
non-adoption. This is correct as far as it goes.

The deeper reading is that the non-Maersk carriers correctly
identified a strategic problem with the platform's governance
shape. A platform controlled by Maersk, with IBM as the technical
operator and Maersk as the principal customer, is structurally a
Maersk asset that other carriers are being asked to provision
with their operational data. The asset's governance accountability
runs to Maersk in any contested case. Whatever the platform's
formal governance arrangements, the controlling-competitor
dynamic is the binding constraint on competitors' participation.

The argument is not that Maersk acted in bad faith or that the
formal governance was inadequate. The argument is that the
platform's shape \emph{is} the problem: a competitor-controlled
infrastructure requires other competitors to ratify the
controller's strategic position relative to the industry, and
there is no formal-governance fix that dissolves this
requirement. MSC and Hapag-Lloyd are not joining a Maersk-controlled
platform regardless of how its governance is configured.

\paragraph{The general lesson.}
TradeLens's failure is not specific to IBM, to Maersk, or to
shipping. It is the recurring pattern of consortium-governance
infrastructure: any platform whose governance accountability runs
to a participant of the industry the platform is meant to serve
will face structural participation refusal from that
participant's competitors. The fix is not better consortium
governance; the fix is no consortium at the platform layer.

Several adjacent industries have run the same experiment with
the same outcome: SWIFT and its alternatives in payments,
various consortium attempts at securities settlement,
healthcare-records consortia. The pattern is robust. The
structural alternative is permissionless, ownerless infrastructure
at the platform layer, with whatever industry-specific arrangements
the participants want to compose on top.

% ════════════════════════════════════════════════════════════════════
\section{The Settlement Primitive}
\label{sec:primitive}

The bonded primitive is developed formally in the companion
mechanism paper, with the deployed Solidity implementation
verified through the companion implementation paper's
multi-method stack (TLA\textsuperscript{+}, symbolic execution,
property-based fuzzing, Certora specifications). We summarize
the features the present argument uses.

\paragraph{Asymmetric bonding (Mechanism~1).}
For a transaction with payment $\pay > 0$ and cumulative upstream
value $\cumval \geq \pay$, the buyer locks $2\pay$ and the seller
locks $2\cumval$. Cooperation is the unique profile surviving
iterated elimination of weakly dominated strategies. The
mechanism scales to $N$-party process trees through progressive
collateralization: each seller bonds against the cumulative value
of all upstream commitments, producing a mesh of independently
secured bilateral edges, each carrying its own equilibrium.

\paragraph{Buyer dominance with atomic resolution (Mechanism~2).}
Only the root buyer can trigger resolution; resolution settles all
active orders in the process simultaneously or not at all. The
all-or-nothing rule induces a weakest-link subgame among sellers,
producing cooperation pressure of magnitude $\pay_i + 2\cumval_i$
on every other seller from each seller's perspective, without
explicit communication, governance, or managerial layer.

\paragraph{What this implies for shipping coordination.}
Mechanism~1 prices each leg's enforcement in proportional bond
posture without requiring any consortium or platform to underwrite
it. Mechanism~2 enforces the multi-leg coordination at resolution
without requiring any party to govern the multi-leg arrangement.
The two mechanisms together supply the coordination function that
TradeLens attempted to supply through a consortium-governed shared
ledger; they supply it without needing the consortium.

% ════════════════════════════════════════════════════════════════════
\section{The Assembly}
\label{sec:assembly}

We now describe the assembly of bonded commitments that
re-architects the supply chain as a permissionless, ownerless
coordination protocol. The assembly has two layers that should be
kept clearly separate: the \emph{organizational DAG} (how the
goods physically flow from factory to consumer, documented via
schema attestations) and the \emph{kernel commit structure} (who
commits to whom on-chain, governed by the kernel's
\texttt{buyer} $=$ \texttt{rootBuyer} rule). The kernel is
DAG-blind; the schema layer is DAG-aware. The two layers co-exist
in one process.

\paragraph{The organizational DAG.}
The supply chain is a sequence of operational handoffs from the
factory to the consumer. We sketch the canonical
international-DTC flow:

\begin{verbatim}
            factory (origin country)
                |
                | handoff: goods packaged, ready to ship
                v
            origin logistics (truck to port)
                |
                | handoff: container stuffed and sealed
                v
            port-of-loading (terminal services)
                |
                | handoff: container loaded
                v
            ocean carrier
                |
                | vessel-position stream (geo-v1),
                | container-seal-v1 sequence
                v
            port-of-discharge (terminal services)
                |
                | handoff: container discharged
                v
            customs broker + customs authority
                |   (broker attests to authority's
                |    determination; authority bonded
                |    for tariff service)
                v
            domestic logistics (warehouse + last-mile)
                |
                | fulfilment-v1 (delivery)
                v
            consumer
\end{verbatim}

Variations on the structure (factory-direct shipment without a
warehouse intermediary, intermodal rail substituted for ocean,
etc.) instantiate the same pattern with the operational divisions
shifted to reflect the route. Every handoff along this DAG is
documented via a schema-typed attestation: \texttt{handoff-v1}
for cargo transfers, \texttt{container-seal-v1} for seal-state
changes, \texttt{geo-v1} for vessel-position updates,
\texttt{jurisdiction-v1} for customs clearance,
\texttt{fulfilment-v1} for delivery confirmation. None of these
attestations is a kernel commit; they are evidentiary records
that document the organizational reality of the goods' flow.

\paragraph{The kernel commit structure.}
The kernel commits, in contrast, are consumer-rooted: the kernel
enforces \texttt{buyer} $=$ \texttt{rootBuyer} in every order in
a process, so every kernel commit in the supply-chain process has
the consumer as buyer. The commits document who the consumer pays
and how much: consumer $\to$ factory, consumer $\to$ origin
logistics, consumer $\to$ origin port authority, consumer $\to$
ocean carrier, consumer $\to$ destination port authority,
consumer $\to$ customs broker, consumer $\to$ customs authority
(for the import tariff), consumer $\to$ domestic logistics,
consumer $\to$ state government (for sales tax), and any other
resource provider the goods draw on. The commits are independent
of who hands the goods to whom in the organizational DAG above.
The consumer alone calls \texttt{resolveProcess} on receipt,
settling the process atomically.

\paragraph{Why this is not a chain of B2B transactions.}
In the conventional supply chain the same goods are sold and
resold at every tier --- factory to wholesaler, wholesaler to
distributor, distributor to retailer, retailer to consumer ---
and each tier extracts margin in exchange for absorbing the
coordination cost between adjacent tiers. The bonded primitive
collapses that coordination cost: the consumer's bonded commits
reach the operational sellers directly, with no margin-extracting
tier in between. The factory is paid by the consumer, not by a
wholesaler. The ocean carrier is paid by the consumer, not by a
shipper-of-record. The firm-shape parties of the conventional
chain (shipper-of-record, distributor, wholesaler, retailer)
become structurally unnecessary at the kernel layer; off-chain
roles such as discovery, curation, brand-as-quality-signal, or
returns handling can still be filled by parties who provide those
specific services, but they cannot insert themselves as kernel
intermediaries between the consumer and the operational sellers.

\paragraph{Per-commit schema clauses.}
Each kernel commit carries the schema clauses relevant to the
seller's role.

\begin{itemize}
  \item \emph{consumer $\to$ factory}: \texttt{commerce-v1}
    payment, \texttt{handoff-v1} on goods packaged for shipment,
    optional \texttt{ghg-protocol-v1} disclosure for the
    manufacturing emissions.
  \item \emph{consumer $\to$ origin logistics}: \texttt{commerce-v1}
    payment, \texttt{handoff-v1} on cargo receipt and on
    container stuffing, \texttt{container-seal-v1} applied at the
    factory dock or origin warehouse.
  \item \emph{consumer $\to$ origin port authority} and
    \emph{consumer $\to$ destination port authority}:
    \texttt{commerce-v1} payment for terminal handling and gate
    access, \texttt{handoff-v1} at port-side transfers, and
    \texttt{container-seal-v1} inspected-intact attestation at
    the port handoff.
  \item \emph{consumer $\to$ ocean carrier}: \texttt{commerce-v1}
    payment, an \texttt{incoterms-2020-v1} clause anchoring the
    shipment term and named place, \texttt{handoff-v1} on cargo
    transfer to the vessel and at discharge,
    \texttt{container-seal-v1} sealing events through the voyage,
    \texttt{geo-v1} vessel-position attestations on the voyage
    cadence, and \texttt{ghg-protocol-v1} disclosure for the
    voyage emissions.
  \item \emph{consumer $\to$ customs broker}: \texttt{commerce-v1}
    payment plus \texttt{jurisdiction-v1} attestation recording
    the customs authority's determination.
  \item \emph{consumer $\to$ customs authority (import duty /
    tariff)}: \texttt{commerce-v1} payment for the import-duty
    obligation. The customs authority is a bonded seller of
    border-services in this commit and a sovereign decisional
    authority on the broker's parallel attestation; the two
    roles are distinct but compose in one process.
  \item \emph{consumer $\to$ domestic logistics}:
    \texttt{commerce-v1} payment plus \texttt{fulfilment-v1}
    attestation at delivery.
  \item \emph{consumer $\to$ state government (sales tax)}:
    \texttt{commerce-v1} payment for the sales-tax obligation.
    The state government is a bonded seller of public services
    funded through the sales tax, on the same composition pattern
    as the customs authority.
\end{itemize}

\paragraph{Operational sub-procurement as separate processes.}
A factory's procurement of materials, labor, and energy; a
carrier's procurement of fuel, crew, and maintenance; a
domestic-logistics provider's procurement of warehouse space and
delivery vehicles --- each of these is a separate process under
the operational party as its own \texttt{rootBuyer}. The kernel
does not link those processes to the consumer's. Cross-process
coordination is achieved through within-process cohort-pressure
on each side and through off-chain reconciliation, not through
cross-process atomic resolution.

% ════════════════════════════════════════════════════════════════════
\section{Bond Posture and Cascade Mechanics}
\label{sec:bonds}

The architecture's structural contribution is that each leg's
risk is priced and propagated by the bonding rule of the kernel,
without consortium intermediation.

\paragraph{Bond posture per commit.}
The kernel maintains a single process-wide cumulative-value
accumulator $G$ that grows monotonically as each commit lands
($G \leftarrow G + P_i$). At every commit the consumer (as
\texttt{rootBuyer}) posts $2P_i$ for that specific payment, and
the seller posts $2G_{\text{at commit time}}$ --- the full
cumulative process value at the moment they commit. Bond posture
is asymmetric: the consumer's per-commit bond scales with the
immediate payment, while each seller's bond scales with the
cumulative value of every upstream commit in the process.
Sellers committing later post correspondingly larger bonds; the
last seller bonds against the full retail value of the goods.

\paragraph{A bond-posture worked example.}
Consider a stylized \$200 imported consumer good (a piece of
consumer electronics manufactured in an exporting country and
delivered to a consumer in an importing country with a 25\%
import duty assessed on the CIF landed cost and a 7.5\% sales
tax assessed on the pre-tax retail price) under one
consumer-rooted process. The table below traces one illustrative
commit sequence.

\begin{table}[h]
\centering
\small
\begin{tabular}{lrrrr}
\toprule
Commit (in sequence) & $P_i$ & $G$ after commit & Consumer bond $2P_i$ & Seller bond $2G$ \\
\midrule
Root: consumer $\rightarrow$ factory (FOB) & $\$85$ & $\$85$ & $\$170$ & $\$170$ \\
Sub: consumer $\rightarrow$ origin logistics & $\$5$ & $\$90$ & $\$10$ & $\$180$ \\
Sub: consumer $\rightarrow$ origin port authority & $\$5$ & $\$95$ & $\$10$ & $\$190$ \\
Sub: consumer $\rightarrow$ ocean carrier & $\$15$ & $\$110$ & $\$30$ & $\$220$ \\
Sub: consumer $\rightarrow$ destination port authority & $\$5$ & $\$115$ & $\$10$ & $\$230$ \\
Sub: consumer $\rightarrow$ customs broker & $\$3$ & $\$118$ & $\$6$ & $\$236$ \\
Sub: consumer $\rightarrow$ customs authority (import duty) & $\$29$ & $\$147$ & $\$58$ & $\$294$ \\
Sub: consumer $\rightarrow$ domestic logistics & $\$39$ & $\$186$ & $\$78$ & $\$372$ \\
Sub: consumer $\rightarrow$ state government (sales tax) & $\$14$ & $\$200$ & $\$28$ & $\$400$ \\
\midrule
Total bond locked across the process & --- & --- & $\$400$ & $\$2{,}292$ \\
\bottomrule
\end{tabular}
\caption{Stylized per-commit bond posture for a \$200 imported
consumer good under one consumer-rooted process. The import duty
($\$29$) is $25\%$ of the CIF landed cost ($\$115 = \$85 + \$5 +
\$5 + \$15 + \$5$, the cumulative process value at point of
arrival at the destination port), not of the retail price. The
sales tax ($\$14$) is $7.5\%$ of the pre-tax retail ($\$186 =
\$200 - \$14$). $G$ is the process-wide cumulative-value
accumulator. The consumer's bond at each commit scales with the
immediate payment to that seller; each seller's bond scales with
the cumulative process value at the moment that seller commits.
The customs authority and the state government are bonded sellers
in their respective commits, on the same composition pattern as
the federal aviation system in the companion scheduled-service
paper. Per-resource payments are illustrative; absolute
magnitudes vary with the goods, route, tariff schedule, and tax
jurisdiction.}
\label{tab:k-bond-posture}
\end{table}

The asymmetry between consumer-side and seller-cohort bonds is
visible: the consumer's aggregate locked bond is $\$400$
($= 2 \times P_{\text{retail}}$, the standard buyer posture),
while the seller cohort's aggregate locked bond is $\$2{,}292$,
nearly six times the consumer's. The seller cohort's bond is a
multiple of the retail value because each seller bonds against
the cumulative $G$ that has grown across upstream commits.

\paragraph{Slip propagation as cohort-compensation pressure.}
When a slip occurs --- a defective unit shipped, a port
congestion delay, a customs hold, a damaged delivery --- the
kernel's atomic-resolution rule means \emph{every} seller's bond
in the consumer's process is at risk: the consumer alone
resolves, and the consumer can withhold resolution while the
seller cohort negotiates compensation. The unit whose performance
failed bears the larger share of cohort-internal compensation;
sellers whose service was unaffected contribute less or not at
all. Compensation flows directly from the seller cohort to the
consumer before any external mechanism (specialized mechanism
contracts, dispute forums, insurance, regulators) is engaged.
The mechanism is the same one developed for scheduled passenger
service in the companion industrial-engineering paper: buyer
dominance plus atomic resolution under bond pressure. The
off-chain question of which operational unit caused the slip ---
factory defect, packaging fault, ocean-leg damage, customs
inspection delay --- is adjudicated by external forums with the
bonded record as input, under the framework of the companion
evidence-and-law paper.

\paragraph{Cohort pressure across the resource provider set.}
The architecture exhibits cohort pressure across the seller set
in the consumer's process. Each seller's bond is at risk if any
other seller's slip propagates to a resolution-withholding by the
consumer. Sellers therefore have direct economic interest in
each other's reliability: a factory whose defect rate forces
repeat compensation negotiations loses repeat consumer business
not because of reputation effects but because the cohort's
expected bond-loss exposure on goods from that factory is direct
and measurable; a port that consistently produces cascading
delays does likewise. The mechanism operates without a
managerial layer; the bonded structure does the work that a
consortium-governed shared ledger was meant to do through
information transparency alone.

% ════════════════════════════════════════════════════════════════════
\section{Schemas}
\label{sec:schemas}

The assembly draws mostly on existing schema infrastructure. Two
new schemas are required: a container-seal attestation and an
Incoterms-2020 anchor. The bill-of-lading transferability
question is treated separately in \Cref{sec:bol}.

\paragraph{Existing schemas.}
The assembly composes the following existing schemas.
\texttt{figaro-commerce-v1} carries the per-leg payment
commitment.
\texttt{figaro-geo-v1} carries pickup / drop-off geohashes and
vessel-position updates.
\texttt{figaro-handoff-v1} carries each container handoff at
transfer points.
\texttt{figaro-jurisdiction-v1} carries the customs-clearance
attestation.
\texttt{figaro-courier-process-v1} carries the carrier's
internal sub-process tree (the multi-stop carrier process).
\texttt{figaro-fulfilment-v1} carries the final delivery to
consignee.
\texttt{figaro-ghg-protocol-v1} (with sister schemas for ISO
14064, PAS 2050, EN 16258) carries per-leg emissions
disclosure on every transport edge.
\texttt{figaro-proximity-policy-v1} together with
\texttt{figaro-proximity-proof-v1} support geofence enforcement
at port handoffs where required.

The composition is non-trivial but draws on established
infrastructure. No agent work is required for these schemas;
they are already on chain and bound to validators.

\paragraph{New schema: \texttt{figaro-container-seal-v1}.}
Container seals are the integrity mechanism for ocean shipping.
A uniquely-numbered seal is applied at the origin and broken only
at customs or at the destination consignee. Multiple parties
(carrier, ports, customs, consignee) need a single canonical
attestation of when the seal was applied, when it was inspected
intact, and if it was breached.

The schema's Layer~A spec carries:

\begin{itemize}
  \item \texttt{containerNumber} (string, ISO 6346 format)
  \item \texttt{sealNumber} (string)
  \item \texttt{event} (enum: applied, inspected\_intact,
    transferred, breached, removed\_by\_customs)
  \item \texttt{inspectorAddress} (address-hex; the wallet
    making the attestation)
  \item \texttt{timestamp} (ISO 8601 datetime)
  \item \texttt{locationGeohash} (string, 8-character precision)
  \item \texttt{evidenceHash} (bytes-hex, optional content hash
    of a photo or inspection report)
\end{itemize}

The schema is the seed of a \emph{container-integrity} family
that could later cover seal-equivalents in other transport modes
(rail container seals, air-freight tamper-evident packaging, etc.)
without reauthoring the family. The schemaId binds permanently
under the \texttt{SchemaRegistry}'s first-write-wins discipline.

\paragraph{New schema: \texttt{figaro-incoterms-2020-v1}.}
Incoterms~2020 (from the International Chamber of Commerce) are
the international trade vocabulary for risk-transfer points and
cost allocation across legs of a shipment: EXW, FCA, CPT, CIP,
DAP, DPU, DDP, FAS, FOB, CFR, CIF
\citep{icc_incoterms_2020, ramberg2010icc}. The vocabulary is
useful as cross-party shared reference, but the traditional
Incoterms semantics come from a context without the bonded
primitive's invariants and may not all map cleanly. We treat
this carefully because it is the canonical worked example of
how traditional commercial vocabulary smuggles in assumptions
the bonded architecture cannot honor.

Three structural mismatches deserve direct treatment.

\emph{Risk transfer is a Figaro-foreign concept.} Incoterms
encode a discretionary risk-transfer point: ``risk transfers from
seller to buyer at point X.'' Under the bonded primitive, risk
follows bond --- the bonded party is at stake until they perform
per the agreed clauses, and there is no separate
\emph{risk-transfer} concept distinct from the bonding rule. A
clean port from Incoterms to bonded clauses requires reinterpreting
``risk transfer at X'' as ``the seller's
\texttt{handoff-v1} attestation due at X discharges the seller's
bond exposure for the upstream segment.''

\emph{Cost-vs-risk separation collapses.} Incoterms separate
``who pays for transport'' from ``who bears risk during
transport.'' Under the bonded primitive, both follow from process
structure: a party who bonds for a sub-process is the buyer of
that sub-process, and risk is what the bonding rule already
prices. Some Incoterms features (CIP/CIF's seller-pays-insurance
feature, for instance) require composition with a parallel
insurance process rather than direct encoding in the
delivery-clause specification.

\emph{Buyer dominance changes the resolution model.} Incoterms
assume good-faith resolution between parties; the bonded
primitive encodes buyer dominance plus mutually-assured
bond-loss as the equilibrium that produces good-faith resolution.
Some Incoterms' implicit dispute-resolution paths --- which
typically rely on local commercial-court adjudication --- continue
to function under the bonded architecture but at the off-chain
forum the parties name in their agreement, not as an internal
property of the term.

The right schema is therefore not ``Incoterm as a traditional
contract clause'' but \emph{Incoterm as a reference to a
Figaro-native delivery-clause specification}. The schema anchors
the term plus the named place; the term-to-delivery-clause
mapping lives in the off-chain ICC publication (the canonical
text) and a runtime term-to-clause table that the assembly's
participants compose from. The validator contract enforces the
anchor (term enum, named place string, place geohash for
cross-reference with \texttt{geo-v1}) but does not encode
per-term semantics in Solidity.

The Layer~A spec carries:

\begin{itemize}
  \item \texttt{term} (enum: EXW, FCA, CPT, CIP, DAP, DPU, DDP,
    FAS, FOB, CFR, CIF)
  \item \texttt{namedPlace} (string; every Incoterm requires a
    named place)
  \item \texttt{placeGeohash} (string, 8-character precision
    for cross-reference with \texttt{figaro-geo-v1})
\end{itemize}

This split lets future Incoterms revisions ship as new schemaIds
(\texttt{figaro-incoterms-2030-v1}, etc.) without mutating prior
anchors. Per-term mapping verification --- which terms map
directly, which require composition with parallel processes
(insurance, customs), and which do not transfer at all --- is
work the schema-author performs against the kernel code before
the schema is registered. Likely outcomes from that verification
are the following.

\emph{EXW, FCA, DAP, DPU, DDP, FAS, FOB} probably map directly,
each becoming a delivery-clause specification: \texttt{handoff-v1}
attestation at the named place, with auxiliary clauses for
customs (DDP) or unloading (DPU) where the term requires them.

\emph{CPT and CFR} map structurally. The ``seller pays carriage''
feature becomes a separate sub-process where the seller is the
buyer of carriage. The ``risk transfers at first carrier''
feature becomes the goods-commit's delivery-clause specification.

\emph{CIP and CIF} map similarly to CPT and CFR plus the
``seller insures for buyer's benefit'' feature. The insurance
assignment may not transfer directly --- it likely requires
composition with a parallel insurance process where the seller
buys insurance from an external insurer naming the buyer as the
beneficiary. The schema-author's verification report identifies
this as a composition requirement rather than as a refusal.

The Incoterms schema is therefore a clean illustration of the
extension-doctrine principle that the protocol anchors shared
references but does not import every contractual assumption that
travels with traditional vocabulary. The vocabulary is useful;
the assumptions are sometimes incompatible; the schema separates
them.

\paragraph{Schemas the assembly does not include.}
The assembly does not include a multi-currency-cross-leg schema
(would break same-unit comparability and the bonding equilibrium),
a force-majeure-escape schema (would re-introduce the kind of
unilateral exit the kernel forbids), or a transferability schema
internal to \texttt{commerce-v1}. The transferability question is
treated separately in \Cref{sec:bol} as a protocol-layer
composition pattern.

% ════════════════════════════════════════════════════════════════════
\section{Bill of Lading Transferability}
\label{sec:bol}

The bill of lading is the carrier's receipt of cargo, the
contract of carriage, and (in the negotiable case) a document of
title. The transferability question --- whether the right to
receive cargo at destination can be transferred mid-shipment from
the original consignee to a new consignee --- is doctrinally
significant because container-shipping commerce often involves
cargo sale in transit, and the canonical legal artifact for that
sale is endorsement of the bill of lading.

The bonded primitive's kernel does not support party substitution
within a single bonded order. Three kernel invariants each
independently rule it out: the buyer is fixed at root commitment
and identical across every sub-order in the process; the parties
to each commitment are fixed at the moment of commit; and the
no-escape-hatches discipline forbids unilateral exit paths from a
bonded state. A literal endorsement of the BoL during shipment
would require any of these invariants to bend, and they do not
bend.

The architectural direction we propose, and which we treat as
\emph{open architectural research} rather than as a deliverable
production-ready pattern, is that BoL transferability is a
\emph{protocol-layer composition} not a kernel-layer property
\citep{bol_research_2026}. The candidate pattern is a
\emph{cancellable-seller wrapper plus counter-process}: the
original buyer, who wants to transfer the right to the cargo
mid-shipment, commits a parallel \emph{cancellation process}
$P_{\text{cancel}}$ in which each sub-seller is compensated for
early exit through a fee owed by the original buyer (and bonded
under the standard kernel rule), and the new buyer commits a
fresh transport process for the remaining legs.

Several open questions hold the pattern back from being a
deliverable mechanism. First, the cancellable-seller wrapper as
sketched in \citet{bol_research_2026} relies on programmatic
acknowledgement of the cancellation by the affected sub-seller
under a pre-agreed fee schedule; whether this constitutes
sub-seller consent in the kernel-discipline sense, or
constitutes a unilateral exit dressed in counterparty clothing,
is unsettled. A sub-seller bonded under a CancellableSeller
wrapper has consented to the wrapper's terms at commit time,
but consent at commit time to a future programmatic
acknowledgement is structurally different from consent at the
moment of cancellation; the design discipline reading depends
on which framing one accepts. Second, the cargo-handoff
operationalization is unaddressed: a fresh transport process
for the remaining legs requires the original carrier to release
cargo to the new carrier, which is a real-world handoff with no
described attestation surface. Third, the cancellation-fee
arithmetic is undefended; setting the fee too low produces an
early-exit option the original sub-sellers will refuse, setting
it too high forecloses the transfer use case, and the
present paper does not defend any specific fee schedule.

The pattern is therefore not a recommended deployment artifact
in the way the schemas of \Cref{sec:schemas} are. We include the
sketch because the BoL-transferability question is the
load-bearing concern for any container-shipping deployment, and
because the conventional answer (negotiable bills of lading
endorsed mid-transit) is precisely what the kernel's
buyer-fixed invariant rules out. The honest position is that
the pattern is open architectural research; deployments
requiring BoL transferability should plan to rely on conventional
BoL doctrine through the off-chain agreement until the pattern
matures, or to constrain their use cases to the
non-transferable region the kernel admits directly.

\paragraph{Why MLETR-style transferable-record analysis applies.}
The companion evidence-and-law paper develops a Class~A /
Class~B records taxonomy that frames the MLETR analysis here
precisely. \emph{Class~A} records are kernel byproducts
(commitment digests, bond deposits, the cumulative-value
accumulator state, resolution events): they are emitted by the
kernel itself, schema-typed at the kernel level, and
cryptographically authenticated. The MLETR Article~10--11
analysis applies directly to Class~A records: kernel structural
properties (singularity through the cumulative-value
accumulator, control through buyer-dominance, integrity through
the chain's consensus) discharge the control and singularity
tests for transferable records. \emph{Class~B} records ---
discretionary schema-validated attestations such as the
\texttt{container-seal-v1} events or
\texttt{handoff-v1} attestations of \Cref{sec:assembly} ---
are evidentiary artifacts \emph{about} the underlying
transferable record (the commitment), not separate transferable
records inheriting MLETR status in their own right. The
cancellable-seller / counter-process pattern preserves the
Class~A fit because the transfer is achieved through bilateral
signing on two separate processes rather than through unilateral
endorsement of a single record; the protocol-layer transfer
pattern does not violate either MLETR test.

\paragraph{What the bonded architecture does and does not preserve.}
A useful organizing distinction is between order-form (named
consignee, transferred by endorsement to a designated party)
and bearer-form (transferred by physical delivery of the
instrument without endorsement) bills of lading. Order-form
endorsement is partially expressible through the cancellable-seller
pattern above (subject to the open questions). Bearer-form BoLs
are functionally equivalent to a
transfer-without-counterparty-consent pattern that the bonded
architecture does not support. Bearer-instrument practices
remain common in commodity trading (oil cargoes, grain shipments)
where the bill is treated as a near-money instrument
\citep{goode_gullifer_2017}; non-bearer order-form practice
predominates in container shipping but is not the only
operational pattern. We treat the bearer-form exclusion as an
honest exclusion: parties whose operations depend on
bearer-instrument liquidity should plan around the conventional
regime for that part of their workflow.

% ════════════════════════════════════════════════════════════════════
\section{Comparison with Existing Approaches}
\label{sec:comparison}

The architecture should be located against the existing landscape
of digital-shipping infrastructures.

\paragraph{TradeLens.}
TradeLens was a permissioned consortium ledger; the architecture
proposed here is a permissionless bonded settlement primitive.
The functional comparison is asymmetric: TradeLens supplied
information transparency without settlement; the bonded
architecture supplies settlement (with information transparency
emerging as a byproduct of the bonded record). The governance
comparison is also asymmetric: TradeLens was governed by
IBM--Maersk; the bonded architecture has no governance at the
kernel layer. The carriers who declined TradeLens did so for a
governance reason; the architecture proposed here removes the
reason.

\paragraph{CargoX.}
CargoX is a blockchain-based bill-of-lading platform that
emphasizes transferable-electronic-record functionality and
MLETR compliance \citep{cargox_2023}. CargoX is a focused product
addressing the BoL transferability problem; the bonded
architecture is a more general settlement primitive whose
relationship to CargoX is composable rather than competitive. A
shipper using the bonded architecture for the bilateral
transport coordination could compose CargoX (or its successor) at
the BoL layer, with the bonded record providing the underlying
settlement evidence and the CargoX layer providing the
transferable-instrument semantics. The two are not in
competition; they address adjacent layers of the stack.

\paragraph{TradeTrust.}
TradeTrust, a Singapore-government-backed framework for
electronic transferable records, supplies open-source tools and
standards for issuing and managing electronic BoLs and similar
instruments \citep{tradetrust_2023}. Like CargoX, it is composable
with the bonded architecture rather than competitive with it.
TradeTrust's standards-based approach to MLETR fit produces
artifacts that are interoperable with the
bonded-architecture's evidence record under the framework
developed in the companion evidence-and-law paper.

\paragraph{MLETR.}
The UNCITRAL Model Law on Electronic Transferable Records (2017)
is a legal framework rather than a technology platform
\citep{uncitral_mletr_2017}. The bonded architecture's fit with
MLETR is treated in the companion evidence-and-law paper. The
short summary is that the kernel-layer artifacts satisfy
Articles~10 and 11 (functional equivalence, control, singularity)
for non-bearer transferable records, with bearer-instrument
semantics noted as an honest exclusion (\Cref{sec:bol}).

\paragraph{The architectural distinction.}
The TradeLens replacement is not a better TradeLens. It is a
re-architecture of the underlying coordination function from
\emph{shared ledger maintained by consortium} to
\emph{bonded settlement primitive composed by participants}. The
shift is structurally substantial: the consortium's role is
absorbed into the kernel's permissionless protocol; the carriers'
participation decision is reduced to whether to compose with the
protocol on individual shipments rather than whether to ratify a
competitor's gatekeeping. The carriers who refused to ratify
TradeLens are not asked to ratify anything in the bonded
architecture; they participate by signing commitments to specific
shipments and decline by not signing. The decision is made at
the level of individual commerce rather than at the level of
industry governance.

% ════════════════════════════════════════════════════════════════════
\section{Scope}
\label{sec:scope}

The argument bounds itself by several scope conditions.

\paragraph{Asset specificity at the port layer.}
Major ports are functional monopolies for the slot, terminal
capacity, and operational infrastructure they provide; carriers
have limited alternative-supplier leverage at congested hubs.
The architecture's competitive properties at the port-supply
layer are constrained accordingly: where the port-of-call is
dictated by the route and the port operator is a near-monopolist,
the buyer-dominance leverage at the sub-procurement is
attenuated. The architecture continues to function in these
cases (the bond posture is well-defined; the resolution path is
clean) but it does not produce price competition in port
services where structural conditions don't admit it.

\paragraph{Customs and sovereign authorities.}
Customs authorities are not Figaro counterparties; their
decisions are sovereign and the architecture treats them as
attestation sources rather than bonded parties. The architecture
does not propose to re-engineer customs administration. It does
provide a verifiable evidence record on which customs decisions
can operate and which customs decisions are recorded against,
which we take as an unambiguous improvement over the contemporary
paper-and-database situation but not as a substitute for the
sovereign authority itself.

\paragraph{Bond-currency at the small-shipper end.}
Small shippers without crypto-treasury infrastructure face
operational friction at the bonding step. The constraint is
real but solvable through standard treasury-of-record
arrangements (a forwarder bonds on behalf of small shippers,
with the small shipper paying conventional fees off-chain to the
forwarder for the bonding service). This is a runtime-tier
implementation question rather than a substrate-tier question;
the substrate is denomination-agnostic.

\paragraph{Insurance and reinsurance.}
Marine insurance, P\&I clubs, and the general apparatus of
shipping-risk underwriting continue to operate. The bonded
architecture changes the risk profile that the marine-insurance
markets price (smaller residual after bond allocation;
concentrated risk at the carrier-bond layer; a new class of
bond-default risk at the port-supplier layer) but does not
displace the insurance function. We expect the marine-insurance
markets to develop fluent products around the bonded
architecture as the architecture diffuses.

\paragraph{Network effects and adoption.}
The architecture's adoption dynamics differ from TradeLens's in
that no participant must ratify any other participant's
governance position. A shipper, forwarder, or carrier can
participate on individual shipments without joining anything;
the participation decision is shipment-by-shipment rather than
platform-membership-shaped. We do not model the diffusion
trajectory beyond noting that the structural barrier that
defeated TradeLens (consortium ratification by competitors) is
absent here.

% ════════════════════════════════════════════════════════════════════
\section{Conclusion}
\label{sec:conclusion}

TradeLens failed for governance reasons that no
consortium-governance fix can resolve, and the structural
alternative is permissionless, ownerless infrastructure at the
platform layer rather than a better consortium. The bonded
settlement primitive of the companion mechanism paper supplies
the architectural shape: each leg of a shipment is a bilateral
bonded commitment; progressive collateralization scales the
primitive across the multi-leg fan-out; buyer dominance plus
atomic resolution propagates damage to the unit causing the
slip; the schemas the assembly composes are largely existing
infrastructure with two well-bounded additions.

The contribution is the assembly DAG, the per-edge mechanism,
the bond-posture mechanics, the schema-set composition, the
Incoterms-as-anchor pattern, and the
cancellable-seller / counter-process pattern for BoL
transferability. None of these is novel mechanism design; the
mechanism-design work was done in the companion mechanism paper.
The contribution is industrial-engineering: a specific
re-architecture of container shipping as a bonded coordination
problem with operational properties demonstrably different from
the consortium-platform regime that has been tried and has
failed.

A natural next step is empirical: a small-corridor pilot
implementation of the architecture for a specific high-volume
trade lane (US-Asia container, intra-EU short-sea, intra-Asia
trade) with measurement of bond-loss frequency, port-supplier
slip rates, and consignee-side recovery proportionality compared
to the conventional regime. We do not undertake the pilot; we
identify it as the experimental work the architecture invites.

\section*{Acknowledgments}
The TradeLens worked example originated in a project repository
of agent-coordination scenarios, with comparative research at
\citet{bol_research_2026}. The present paper formalizes the
scenario into an industrial-engineering register.

% ════════════════════════════════════════════════════════════════════
\begin{thebibliography}{99}

\bibitem[Jensen et al.(2019)]{jensen_etal_2019}
Jensen, T., Hedman, J., and Henningsson, S.
\newblock How TradeLens Delivers Business Value With Blockchain
  Technology.
\newblock \emph{MIS Quarterly Executive}, 18(4):221--243, 2019.

\bibitem[Sarker et al.(2021)]{sarker_etal_2021}
Sarker, S., Henningsson, S., Jensen, T., and Hedman, J.
\newblock Use Cases for Blockchain in the Enterprise: TradeLens.
\newblock \emph{MIS Quarterly Executive}, 20(4):257--276, 2021.

\bibitem[Lloyd's Loadstar(2022)]{lloyd_loadstar_2022}
The Loadstar.
\newblock IBM and Maersk to Sink TradeLens Blockchain Platform.
\newblock The Loadstar, November~30, 2022.

\bibitem[ICC(2020)]{icc_incoterms_2020}
International Chamber of Commerce.
\newblock \emph{Incoterms 2020: ICC Rules for the Use of Domestic and
  International Trade Terms}.
\newblock ICC Publication 723E, 2020.

\bibitem[Ramberg(2010)]{ramberg2010icc}
Ramberg, J.
\newblock \emph{ICC Guide to Incoterms 2010}.
\newblock ICC Publication 720E, 2010.

\bibitem[CargoX(2023)]{cargox_2023}
CargoX d.o.o.
\newblock \emph{CargoX Platform: Blockchain Document Transfer
  Solution for Shipping}.
\newblock CargoX product documentation, 2023.

\bibitem[TradeTrust(2023)]{tradetrust_2023}
Government of Singapore.
\newblock \emph{TradeTrust: Open-Source Framework for Electronic
  Transferable Records}.
\newblock TradeTrust technical documentation, 2023.

\bibitem[UNCITRAL(2017)]{uncitral_mletr_2017}
United Nations Commission on International Trade Law.
\newblock \emph{UNCITRAL Model Law on Electronic Transferable
  Records}.
\newblock UN Publication, 2017.

\bibitem[BoL Research(2026)]{bol_research_2026}
Daliana, A.
\newblock Bill of Lading Research: CargoX, TradeTrust, MLETR,
  TradeLens vs Figaro.
\newblock Internal research deliverable, April 28, 2026.

\bibitem[Goode and Gullifer(2017)]{goode_gullifer_2017}
Goode, R. and Gullifer, L.
\newblock \emph{Goode and Gullifer on Legal Problems of Credit and
  Security}, 6th edition.
\newblock Sweet \& Maxwell, London, 2017.

\end{thebibliography}

\end{document}
